% !PRINTED 241
\weberpage{241}
3. Es ist
\[
F_n[j(n\omega), j(\omega)] = 0
\]
und wenn wir $n\omega$ durch $\omega$ ersetzen:
\[
F_n\Bigl[j(\omega), j\Bigl(\frac{\omega}{n}\Bigr)\Bigr] = 0.
\]
Da nun $j\left(\dfrac{\omega}{n}\right)$ gleichfalls eine Wurzel der Invariantengleichung ist, so folgt, daß die beiden Gleichungen $\nu$ten Grades
\[
F_n[x, j(\omega)] = 0, \quad F_n[j(\omega), x] = 0
\]
eine Wurzel gemeinsam haben, und folglich, wegen der Irreducibilität der ersteren und der Gleichheit des Grades, alle. Es ist daher
\[
F_n(x, y) = C F_n(y, x)
\]
für unbestimmte Werte der Variablen $x$, $y$ und ein konstantes $C$.
Daher auch, durch Vertauschung von $x$ mit $y$:
\[
F_n(y, x) = C F_n(x, y),
\]
also $C^2 = 1$ oder
\[
F_n(x, y) = \pm F_n(y, x).
\]
Wenn wir nun $x = y$ setzen, so folgt:
\[
F_n(y, y) = \pm F_n(y, y),
\]
also, wenn das untere Zeichen gilt,
\[
F_n(y, y) = 0.
\]
Dies ist aber nur möglich, wenn $n = 1$ ist [wo $F_1(x, y) = x - y$ wird], da sonst $F_n(x, y)$ durch $x - y$ teilbar sein müßte, was der Irreducibilität widerspricht. Daher ist immer, sobald $n > 1$ ist:
\weq{18}{F_n(x, y) = F_n(y, x).}

4. Es sei
\[
n = n' n'', \quad \nu' = \psi(n'), \quad \nu'' = \psi(n''),
\]
und $n'$, $n''$ \gesperrt{ohne gemeinsamen Teiler}, und $x_1, x_2, \dots x_{\nu''}$ seien die Wurzeln der Gleichung
\weq{19}{F_{n''}[x, j(\omega)] = 0.}
Das Produkt
\weq{20}{F_{n'}(x, x_1)\, F_{n'}(x, x_2) \dots F_{n'}(x, x_{\nu''}),}
dessen Grad in bezug auf $x$ gleich
\[
\nu = \psi(n) = \psi(n') \psi(n'')
\]
ist, hängt als symmetrische Funktion der Wurzeln von (19), rational von $j(\omega)$ ab und verschwindet für
\[
x = j\left(\frac{\omega}{n' n''}\right),
\]
%CONT>>
